% 本文件由 LaTeX 排版 Agent 根据有效 translation.json 自动生成。
% author=Gregory Thaumaturgus; work_id=0609; title=四篇讲道词

\chapter{四篇讲道词}

% structure: p0001 source=h1 target=omitted reason=page-title-already-rendered-by-parent

\Emph{讲道词 1} 论圣母领报\newline{}\Emph{讲道词 2} 论圣母领报\newline{}\Emph{讲道词 3} 论圣母领报\newline{}\Emph{讲道词 4} 论主受洗\par

\SourceRecord{Translated by S.D.F. Salmond.；Ante-Nicene Fathers；Vol. 6.；Edited by Alexander Roberts, James Donaldson, and A. Cleveland Coxe.；Buffalo, NY: Christian Literature Publishing Co.,；1886.；<http://www.newadvent.org/fathers/0609.htm>.}

% structure: p0001 source=h1 target=section reason=page-title

\section{第一篇讲道}

论天神报喜于圣童贞圣母玛利亚\par

今日，天使的唱诗班欢唱赞美的旋律，基督降临的光明照耀在信友身上。今日，对我们而言是欢乐的春天，正义的太阳基督以明净之光普照我们，照亮了信友的心智。今日，亚当被再造，在天使的唱诗班中行走，展翅飞向天堂。今日，整个大地充满了喜乐，因为圣神的临在已经实现于人间。今日，天主的恩宠和不可见之事的望德，超越一切想象，透过所有奇迹闪耀，使我们清楚辨明那从永恒隐藏的奥秘。今日，永不凋谢的美德花环被编织。今日，天主愿意为那些乐于聆听祂、喜爱祂节庆的人戴上圣冠，召叫热爱坚定不移信德的人作为祂的蒙召者和继承者；天国急切地召唤那些思念天上之事的人，加入无形唱诗班的神圣事奉。今日，达味的话应验了：\BibleQuote{愿诸天欢乐，愿大地踊跃！愿田野欢喜，愿森林中的一切树木在上主面前欢乐，因为祂来了。}达味提及树木；主的先驱也提到它们，说它们\BibleQuote{应当结出与悔改相称的果实}，或者更确切地说，为主来临的果实。但我们的主耶稣基督应许所有信祂的人永远的喜乐。因为祂说：\BibleQuote{我要再见你们，你们要喜乐，并且你们的喜乐没有人能夺去。}今日，基督徒那荣耀而不可言喻的奥秘被明确地向我们启示，他们甘愿将希望如同印玺一样放在基督身上。今日，那站在天主面前的加俾额尔来到纯洁的童贞女那里，向她传达喜讯：\BibleQuote{万福！充满恩宠者！她心中思忖这是怎样的请安。天使立刻接着说：主与你同在。不要害怕，玛利亚，因为你已在天主前蒙恩。看，你将怀孕生子，给祂起名叫耶稣。祂将是伟大的，被称为至高者的儿子；上主天主将把祂祖先达味的王位赐给祂，祂将统治雅各伯家直到永远：祂的王权没有终结。玛利亚对天使说：这事如何成就？因为我不认识男人。}我仍将保持童贞吗？我的童贞荣誉不会因此丧失吗？当她还在为这些事困惑时，天使简短地将整个信息的要点放在她面前，对纯洁的童贞女说：\BibleQuote{圣神要降临到你身上，至高者的能力要荫庇你，因此那要出生的圣者，将被称为天主子。}因为是什么，就必将被称为是什么。于是，恩宠谦和地从所有世代中独选了纯洁的玛利亚。因为她在一切事上确实显出明智；所有世代中没有女人曾像她一样出生。她不像原始的童贞女厄娃，在乐园独处时，心不在焉，轻率地听从了万恶之源——蛇的话，从而在心思上堕落；通过她，那欺骗者释放毒液，拒绝与之同死，将毒液带入整个世界；圣人们的一切困苦皆由此而来。但在圣童贞女身上，那（第一母亲）的堕落得到了修复。然而，这位圣者竟不能领受这恩赐，除非她先知道是谁派遣它，这恩赐是什么，以及谁传递它。当圣者困惑地思量这些事时，她对天使说：\Quote{你从哪里给我们带来这样的祝福？从哪座宝库里发出了言语的珍珠？这恩赐从哪里获得对我们的旨意？你从天而来，却行走在地上！你显现人的形像，却闪耀着耀眼的光芒。}圣者心中思量这些事，总领天使解决了这种推理中的困难，对她说：\Quote{圣神要降临到你身上，至高者的能力要荫庇你。因此那要出生的圣者，将被称为天主子。不要害怕，玛利亚；因为我来不是要惊吓你，而是要驱除恐惧。不要害怕，玛利亚，因为你已在天主前蒙恩。不要以自然的标准来质询恩宠。因为恩宠不愿受自然法则的约束。你知道，玛利亚，那些向先祖和先知隐藏的事。你知道了，童贞女，那些至今向天使隐藏的事。你听到了，最纯洁者，连受默感者的唱诗班都不曾配得的事。梅瑟、达味、依撒意亚、达尼尔和所有先知都预言了祂，但方式他们不知道。然而你，最纯洁的童贞女，如今独享了所有这些人所不知道的事，并知道了它们的起源。因为哪里有圣神，那里一切就都有条不紊。哪里有神圣恩宠，那里一切在天主都是可能的。圣神要降临到你身上，至高者的能力要荫庇你。因此那要出生的圣者，将被称为天主子。}如果祂是天主子，那么祂也是天主，与圣父同一形态，\OriginalTerm{同永恒}{co-eternal}；在祂内，圣父拥有一切彰显；祂是圣父的位格肖像，通过祂的映照，（圣父的）荣耀闪耀。如同从永不枯竭的泉源流出溪流，同样，世界的光，永恒而真实的光，即我们的天主基督，也从这永不枯竭、永活的泉源流出。因为先知们曾宣讲：\BibleQuote{河流的支流使天主的城邑欢乐。}并不仅是一座城，而是所有的城；因为正如它使一座城欢乐，同样也使整个世界欢乐。因此，天使首先对童贞玛利亚说：\BibleQuote{万福！充满恩宠者，主与你同在；}因为在她那里珍藏了恩宠的全部宝藏。因为所有世代中，只有她作为童贞女，身心灵纯洁地兴起；只有她怀着那以言语承载万有的那一位。这位圣者的美丽不仅在于身体令人赞叹，还有灵魂的内在德行。因此，天使先用这问候向她致意：\Quote{万福！充满恩宠者，主与你同在，没有地上的配偶；}祂自己与你同在，祂是圣化的主，纯洁之父，不朽的创始者，自由的施予者，救恩的监护者，真正平安的管理者和供应者，祂用处女地造人，用人的肋骨造厄娃。这位主与你同在，另一方面也出于你。因此，亲爱的弟兄们，让我们接过天使的颂歌，尽力回报应得的赞美，说：\BibleQuote{万福！充满恩宠者，主与你同在！}因为你实在应当喜乐，因为天主的恩宠，如祂所知道的，选择与你同在——荣耀之主与婢女同在；\Quote{祂比世人更美，\InnerQuote{与美丽的\Emph{童贞女}。}祂圣化一切，与无玷者同在。天主与你同在，与你同在的还有那完全的人，在他内住着整个天主性的圆满。万福！充满恩宠者，那照亮所有信祂之人的光明之泉！万福！充满恩宠者，那理性之日的升起，生命的不朽花朵！万福！充满恩宠者，那馨香的草地！万福！充满恩宠者，那常青的葡萄树，使敬爱你的灵魂喜乐？万福！充满恩宠者！那未经耕种却结出丰硕果实的土壤：因为你按自然规律生育，如同我们一样，按定期的过程，却又超越自然，或说高于自然，因为天主圣言从上降临到你内，在你神圣的胎中形成新亚当，并且由于圣神赐予圣童贞女怀孕的能力；祂的身体的真实性取自她的身体。正如珍珠来自两种本性，即闪电和水，海洋的隐秘标记；同样，我们的主耶稣基督来自纯洁、贞洁、无玷、圣洁的童贞玛利亚，无混合、无变化：完全的天主性和完全的人性，在一切方面与圣父相等，在一切方面与我们\OriginalTerm{同体}{consubstantial}，除了罪。}\par

多数圣教父、圣祖和先知渴望看见祂，并作祂的目击者，但没有达到。他们中的一些人在异象中在\OriginalTerm{预象}{type}和模糊中看见祂；又有些人，得以通过云彩媒介听到天主的声音，并蒙恩看见圣天使的景象；但惟有纯贞的\Person{玛利亚}，总领天使\Person{加俾额尔}光辉地显现自己，给她带来喜乐的问候：\BibleQuote{万福，满被圣宠者！}于是她接受了这话，并在依照身体过程的实现之时，生下了无价珍珠。那么，亲爱的诸位，你们也来吧，让我们歌唱那由受默感的达味竖琴教给我们的旋律，并说：\BibleQuote{上主，请起身，到你的安息之所；你和你圣所的约柜。}因为童贞玛利亚实在是一个\OriginalTerm{约柜}{ark}，内外包金，收纳了\OriginalTerm{圣所}{sanctuary}的整个宝藏。\BibleQuote{上主，请起身，到你的安息之所。}上主，请从圣父的怀抱中起来，为的是要举起那首先受造的人堕落的种族。达味在预言中将这些事陈明，对那要从他自己而出的\OriginalTerm{枝条}{rod}，并要发芽成为那美丽果实的\OriginalTerm{花朵}{flower}的，说：\BibleQuote{女儿，请听，请看，请侧耳，忘记你的民族和你的父家；这样，君王将爱慕你的美丽：因为祂是你的主天主，你当敬拜祂。}女儿，请听那从前关于你预言的事，为叫你也用理解的双眼看到这些事本身。请听我向你预告的事，也请听那总领天使向你明确宣告的完美奥秘。那么，亲爱的诸位，让我们回想以前的事；并赞美、庆祝、颂扬、祝福那从\Person{叶瑟}奇妙地生出的枝条。因为\Person{路加}在受默感的福音叙述中，不仅为\Person{若瑟}作证，也为天主之母\Person{玛利亚}作证，并论及达味的家族和家系，如此记载：他说：\BibleQuote{若瑟从\Place{加里肋亚}，到了\Place{犹太}的一座名叫\Place{白冷}的城，要同他聘定的、怀孕的玛利亚去登记，因为他们都是达味家族和家系的人。于是，在那里的时候，她分娩的日子满足了；她生了一个儿子，是所有受造物的首生者，用襁褓包起来，放在马槽里，因为在客栈中没有地方给他们。}她用襁褓包裹那以光为衣的祂。她用襁褓包裹那创造万有的祂。她把那坐在\OriginalTerm{革鲁宾}{cherubim}之上、受万千天使赞美的祂放在马槽里。在为哑畜预备的马槽里，\OriginalTerm{天主的圣言}{Word of God}安息，为要把真理性的知觉赐给那些藉自由意志实在不理性的人们。在牲畜吃食的板上，放置了\OriginalTerm{天上的食粮}{heavenly Bread}，为给那些像地上走兽一样生活的人们提供对精神养分的分享。在客栈中甚至连地方也没有给祂。祂没有找到地方，祂曾以祂的话建立了天地；\BibleQuote{因为祂本是富有的，却为了我们而成为贫穷的}，并为我们本性的救恩选择了极度的卑下，出于祂对我们固有的良善。那在圣父怀中、在天上完成了不可言喻的\OriginalTerm{救恩计划}{economy}奥秘的整个安排，并在山洞里、在母亲怀抱中的祂，安息在马槽里。天使的歌咏团环绕祂，歌唱在天上的光荣和地上的和平。在天上，祂坐在圣父的右边；在马槽里，祂安息，就如同在革鲁宾之上。那里实在是祂革鲁宾的宝座；那里是祂的君位。至圣所，地上唯一荣耀的，比圣所更圣的，是我们的天主基督安息的地方。愿光荣、尊威、权能归于祂，同无玷的圣父，及至圣并赋予生命的圣神，自今至永远，及世之世。阿们。\par

\SourceRecord{Translated by S.D.F. Salmond.；Ante-Nicene Fathers；Vol. 6.；Edited by Alexander Roberts, James Donaldson, and A. Cleveland Coxe.；Buffalo, NY: Christian Literature Publishing Co.,；1886.；<http://www.newadvent.org/fathers/06091.htm>.}

% structure: p0001 source=h1 target=section reason=page-title

\section{第二篇讲道}

论向童贞玛利亚报喜\par

我们有义务将一切庆节和赞美颂，如同祭献一般呈献给天主；而首先，是向天主之母的报喜，即天使向她所作的问候：\BibleQuote{万福！满被恩宠者！} 因为在新约中，一切智慧与救恩教理的首要便是这问候：\BibleQuote{万福！满被恩宠者！}是从光之父那里传给我们的。而这句称呼\Quote{\InnerQuote{满被恩宠}}涵盖了全人类的本性。\BibleQuote{万福！满被恩宠者}，在神圣的受孕和荣耀的怀胎中，\BibleQuote{我给你们报告一个为全民族的大喜讯}。而主，祂来是为了完成救赎的苦难，曾说：\BibleQuote{我要再见你们，你们心里就喜乐了，并且你们的喜乐谁也不能从你们夺去}。而在祂复活后，又藉圣妇们的手，首先给了我们\Quote{万福！}的问候。再者，宗徒也以相似的话宣报说：\BibleQuote{应常欢乐，不断祈祷，事事感谢}。所以，亲爱的诸位，请看主是如何处处不可分割地赐予我们那超乎想象、永恒不绝的喜乐。因为既然圣童贞在肉身生活中拥有不朽的公民身份，并如此在一切德行中行事，且过着超越常人标准的生活；因此，那出自天主圣父的圣言认为适合取肉体，并从她身上取得完美的人性，为的是在罪藉以进入世界的同一肉体中——死亡亦由罪而来——罪能在肉体中被定罪，并且罪的诱惑者能在圣体的埋葬中被战胜，同时也藉此彰显复活的开始，并在世上建立永生，为人与天主圣父建立共融。我们在此要说什么？要略过什么？谁能解释这奥迹中不可理解的部分？但目前让我们回到主题。加俾额尔被派遣到圣童贞那里；无形体者被差遣到那位在肉体中追求不朽生活、在纯洁和德行中度日者那里。当他来到她面前时，他首先以问候对她说：\BibleQuote{万福！满被恩宠者！主与你同在。} 万福！满被恩宠者！因为你确实做了堪当喜乐的事，因为你穿上了纯洁的外衣，腰束谨慎的带子。万福！满被恩宠者！因为你竟成了天上喜乐的载体。万福！满被恩宠者！因为藉着你，喜乐被宣告给整个受造界，人类也藉着你恢复了原始的尊严。万福！满被恩宠者！因为万物的创造主要被怀抱在你的手臂中。她因这话感到困惑，因为她不熟悉人们的各种言辞，并喜爱安静，作为谨慎与纯洁之母；但她自己既是纯洁、无玷、无瑕的肖像，并不像大多数先知那样因天使的出现而惊恐退缩，因为真正的童贞确实与天使有某种亲和与平等。因为圣童贞小心守护着童贞的火炬，并殷勤留意，不使它熄灭或被玷污。正如一个穿着灿烂长袍的人认为不让任何不洁或污秽碰触它是一件大事，同样，圣玛利亚也思忖着说：\Quote{这番关注意味着什么深远的图谋或诱惑的意图？难道这\InnerQuote{万福}一词会像古时蛇魔给厄娃的\InnerQuote{如同天主一样}的美妙许诺那样，成为我烦恼的起因吗？难道那万恶之源的魔鬼又变成了光明的天使，因忌恨我未婚夫杰出的节制，以貌似亲切的言辞攻击他，却发现自己无力战胜如此坚定的心灵、无法欺骗那人，于是转而攻击我，以为我是更容易受影响的心灵？难道这\InnerQuote{万福}（恩宠与你同在）一词，是将来无恩宠的记号？这祝福与问候是以讽刺方式说出的吗？难道蜂蜜里没有隐藏着毒药？难道这不是报喜者的言辞，而其结局要使成为设计者的猎物？他怎么能够向一个他不认识的人如此问候呢？}她困惑地思量这些事，并说了出来。那时，总领天使又用那所有人都可相信、永不消逝的喜乐宣告对她说话，对她说：\BibleQuote{玛利亚，不要怕，因为你在天主前获得了宠幸。} 你立刻有了所说之事的证明。因为我不仅让你明白无可恐惧，还向你展示了消除一切恐惧原因的钥匙。因为藉着我，所有天上的掌权者都向你——圣童贞——致敬；不，是祂自己，一切天上掌权者和所有受造物的主，拣选了你为圣者和全然美丽者；并且藉着你神圣、贞洁、纯洁、无玷的子宫，那光照的珍珠为普世救恩而出世：因为在全人类中，你生来就是圣人，比任何其他人都更尊贵、更纯洁、更虔诚；你有比雪更白的头脑，有比任何精金更纯净的身体，还有如同厄则克耳所见并描述的那样的子宫：\BibleQuote{在活物头顶上，有像穹苍的样式，像耀眼的晶体的样子，铺张在它们头顶上。在穹苍之上，有一种像宝座的样式，像蓝宝石的样子；在宝座之上，有一个像人的样式。我望见一种像琥珀的形色，在宝座四周有像火的形状。}\par

显然，先知以预像看见了那由圣童贞所生者。啊，圣童贞，若非你那时以一切荣耀与美德发出光辉，你本无力承受祂。那么，我们该用什么赞美之词来描述她的童贞尊严？用什么赞美之标和宣扬来颂扬她的无玷形象？用什么属灵诗歌或言语来尊崇她在天使中最光荣的那一位？她如同圣神所荫庇的结实橄榄树，栽种在天主的殿中；藉着她，我们被称为基督国度的儿子和继承人。她是不朽的永春乐园，园中栽种了生命树，为众人提供永生的果实。她是童贞女的夸耀和光荣，是母亲们的喜乐。她是信者的坚固依靠，虔诚者的扶助。她是光明之衣，德行之居。她是永流不息的泉源，生命之水从中涌出，产生了主降生成人的显现。她是正义的纪念碑；凡爱慕她、倾心于童贞般的纯真和纯洁的人，必享天使的恩宠。凡戒绝酒醉和烈酒之纵乐的人，必因生命之树的产物而喜乐。凡保持童贞之灯不灭的人，必得接受不朽的永生不凋花冠。凡拥有节制之无玷衣袍的人，必被迎入正义的奥秘新房。凡比他人更接近天使等级的人，也必更真实地享受主的真福。凡拥有理解之明灯和良心之纯净香品的人，必继承属灵恩宠和属灵嗣子的应许。凡相称地庆祝天主之母童贞玛利亚领报瞻礼的人，必获得相称的酬报，就是更圆满地分享那句问候：\BibleQuote{万福！充满恩宠者！}因此，我们有责任庆祝这个庆节，因为它以喜乐和欢愉充满了整个世界。让我们以圣咏、赞美诗和属灵歌曲来庆祝。昔日以色列人也庆祝他们的庆节，但那时用无酵饼和苦菜，关于此，先知说：\BibleQuote{我要将你们的庆节变为悲哀和哀号，将你们的喜乐变为羞耻。}但我们的主向我们保证，祂要藉着痛悔的果实，将我们的悲哀变为喜乐。再者，第一个盟约维持了敬神礼仪的正义要求，如同在先祖亚巴郎身上那样；但这些要求在于割礼带来的肉体痛苦，直到圆满之时。\BibleQuote{法律是藉着梅瑟赐给他们的}，为管教他们；\BibleQuote{但恩宠和真理}却是由耶稣基督赐给我们。这一切福泽的开始，出现在对充满恩宠者玛利亚的报喜中，出现在堪受一切赞美的救主的\OriginalTerm{救恩计划（或经世）}{economy}中，以及祂神圣且超尘世的训导中。从那里，理解之光的光线照在我们身上。从那里，为我们涌出智慧和永生的果实，发出清晰纯净的虔诚溪流。从那里，来到我们中间的，是神圣知识宝库的光辉灿烂。\BibleQuote{因为永生就是：认识你，唯一的真天主，和你所派遣来的耶稣基督。}又说：\BibleQuote{你们查考圣经，因为你们认为里面有永生。}正是为此，天主知识的宝库向查考神圣神谕的人启示。那受默感的圣经的宝库，护慰者今日为我们开启了。愿先知的口舌和宗徒的训导成为我们智慧的宝库；因为没有法律和先知，或没有福音作者和宗徒，就不可能有救恩的确切希望。因为我们的主藉着圣先知和宗徒的口舌说话，天主喜爱圣人们的话语；并非祂需要人的言语，而是祂喜悦善好的心志；并非祂从人获得什么益处，而是祂在义人正直的灵魂中找到安息的满足。因为基督并非因我们所说的而显为伟大；而是我们因蒙受祂的恩惠，以感恩之心宣扬祂对我们的恩泽；并非我们能达到其中相称的地步，而是我们尽己所能给予应得的回报。因此，当福音或书信被宣读时，不要让你们的注意力集中于书本或读经员，而应集中于从天上向你们说话的天主。因为书本是可见的，而基督是所说的神圣对象。因此，它给我们带来了那堪受一切赞美的救主的\OriginalTerm{救恩计划（或经世）}{economy}的喜讯，即：祂虽是天主，却因对人类的慈爱而成为人，并没有舍弃祂从永远所拥有的尊严，而是承担了那施行救恩的计划。它给我们带来了那堪受一切赞美的救主的\OriginalTerm{救恩计划（或经世）}{economy}的喜讯，即：祂如同一位为病患而来的医生，不是用药剂医治他们，而是藉着祂仁爱的倾注恢复了他们。\par

它给我们带来了这全然可赞的救主救恩计划的喜讯，即：给那些迷失方向的人指出了救恩之路，并向绝望者宣告了救恩的恩宠，这恩宠以不同方式祝福众人：寻找迷失者，光照盲者，给死者生命，释放奴隶，赎回俘虏，成为我们众人的一切，好作我们真正的救恩之道。祂所做这一切，并非因我们对祂的善意，而是出于施惠者自身固有的慈爱。因为救主所做的一切，并非为了自己获取德行，而是为了使我们获得永生。祂确实按照天主的肖像造了人，并命其居住在乐园中。但人受魔鬼欺骗，成了天主旨意的违犯者，遂被判处死刑。因此，他的后代由于继承关系，也分担了父亲的罪责，并被要求承担定罪的账目。\BibleQuote{因为死从亚当直到梅瑟掌权了}。但主以祂对人的慈爱，当祂看到自己所造的受造物被死亡权势所控制时，并未最终离弃祂按自己肖像所造的人，反而在每一世代探望他，未曾抛弃他；祂首先在圣祖中显现自己，然后在法律中宣告自己，在先知中呈现自己的肖像，预表了救恩的计划。当时期一满，祂荣耀显现的时候到了，祂预先派遣总领天使加俾额尔向童贞玛利亚报喜。他从上面不可言喻的权势下降到圣童贞面前，首先以问候语对她说：\BibleQuote{万福！充满恩宠者}。当这句\BibleQuote{万福！充满恩宠者}传到她耳中时，就在她听到的瞬间，圣神进入了童贞无玷的圣殿，她的理智和肢体一同被圣化。自然界站在对面，自然的结合退居一旁，惊愕地注视着自然的主，以违反自然、不如说超乎自然的方式，在身体中行奇妙之事；基督正是用魔鬼用来攻击我们的同一武器拯救了我们，祂取了我们\OriginalTerm{可受苦的}{passible}的身体，为将那更大的恩宠赐予那缺少它的受造物。\BibleQuote{罪恶在哪里越多，恩宠在那里也越格外丰富}。恩宠恰当地被赐予了圣童贞。因为这在福音史的神谕中也记载着：\BibleQuote{到了第六个月，天使加俾额尔奉天主差遣，往加里肋亚一座名叫纳匝肋的城去，到一位童贞女那里，她已与达味家族中的一个名叫若瑟的男子订了婚，童贞女的名字叫玛利亚}等等。而这对于圣童贞来说是第一个月。正如圣经在法律书中说：\BibleQuote{你们要以本月为正月，为一年之首}。\BibleQuote{你们世世代代要向主举行逾越节庆节}。它也是匝加利亚的第六个月。所以，圣童贞果然出自达味家族，家住白冷，并按照亲属关系合法地与若瑟订了婚。她的未婚夫是她的守护者，也拥有她无玷的纯真。而圣童贞的名字非常相称。因为她被称为玛利亚，按解释就是\Emph{\OriginalTerm{光照}{illumination}}。有什么比童贞之光更明亮呢？为此，那些力求正确把握德行真义的人也称德行为童贞。但如果拥有童贞之心是这么大的恩赐，那么灵魂连同那珍视童贞的身体将是多么大的福分啊！\par

这样，荣福童贞玛利亚虽仍在肉躯之内，却保持了不朽的生活，并以信德接受了总领天使所传报的事。此后，她殷勤地往山区去见她的亲戚依撒伯尔。\\BibleQuote{她进了匝加利亚的家，就给依撒伯尔请安}，仿效了天使。\\BibleQuote{一听见玛利亚请安，依撒伯尔遂觉得腹中的婴儿欢喜跳跃，且充满了圣神。} 玛利亚的声音带着权能行事，使依撒伯尔充满圣神。她的舌头仿佛一口永不枯竭的泉源，向亲戚涌流出预言性的恩赐之流；当时腹中的胎儿还被束缚在子宫里，却已预备要舞蹈跳跃。那是奇妙欢跃的标记。因为无论蒙恩的女子身在何处，她都以喜乐充满一切。\\BibleQuote{依撒伯尔大声呼喊说：\\InnerQuote{在女人中你是蒙祝福的，你的胎儿也是蒙祝福的。吾主的母亲驾临我这里，这是从那里来的呢？在女人中你是蒙祝福的。} }因为你成了女人中新的受造的开始。你给了我们进入乐园的坦然无惧，驱散了我们古老的愁苦。因为在你以后，女人的族类不再受责难。夏娃的后嗣不再惧怕古老的诅咒或分娩的痛苦。因为基督，我们族类的救赎者，万性的救主，属灵的亚当，祂医治了地上受造物的创伤，从你神圣的胎中出生。\\BibleQuote{在女人中你是蒙祝福的，你的胎儿也是蒙祝福的。} 因为为我们带来万福的那一位，显明是你所结的果实。我们从那位不育者的清晰话语中读到这些，而荣福童贞本人更清晰地表达了这一点，当她向上主呈上充满感恩、接纳和神知的颂歌时，她将古老的事与新生的事一同宣告，将往昔的事与万世终结的事一同传扬，并在短短的一篇话中总结了基督的奥秘。\\BibleQuote{玛利亚说：\\InnerQuote{我的灵魂颂扬上主，我的心神欢跃于天主，我的救主，}} 等等。\\BibleQuote{祂扶助了祂的仆人以色列，纪念祂的仁慈，以及对亚巴郎和他的子孙永远的盟约。} 你看到荣福童贞如何超越了先祖的成全，她如何证实了天主与亚巴郎所立的盟约，当时祂说：\\BibleQuote{这是我与你之间所立的盟约。} 为此，祂来并证实了与亚巴郎的盟约，在祂自己身上奥秘地接受了割损的标记，并显明自己是法律和先知的满全。因此，这首预言之歌，天主之母向天主呈上，说：\\BibleQuote{我的灵魂颂扬上主，我的心神欢跃于天主，我的救主：因为全能者在我身上行了大事，祂的名字是圣的。} 因为祂使我成为天主之母，也保全了我做童贞女；藉着我的子宫，万代的圆满一同被收纳以得圣化。因为祂祝福了每个时代，不论男女、青年、少年、老人。\\BibleQuote{祂用自己的手臂施展大能，} 为我们对抗死亡和魔鬼，撕毁了我们罪债的字据。\\BibleQuote{祂把心高气傲的人驱散，} 诚然，祂驱散了魔鬼本身及所有服事他的恶魔。因为他在心里极其狂傲，竟敢说：\\BibleQuote{我要把我的宝座放在云上，我要与至上者相似。} 至于祂如何驱散了他，先知在接下来的话中指明了，说：\\BibleQuote{但是你却被抛下阴间，} 连同你所有的随从。因为祂处处推翻了他们的祭坛和虚幻之神的敬拜，并为祂自己从外邦民族中预备了一个特选的子民。\\BibleQuote{他从高位上推下有权势者，却举扬了卑微之人。} 这些话简要地暗示了犹太人的被弃和外邦人的接纳。因为犹太人的长老和经师，以及享有其他特权的人，因滥用财富和无法无天使用权力，被祂从每一个座位上——无论是预言的、司祭的、立法的或教义的——被推下，并被剥夺了所有祖先的财富、祭献、众多的庆节及王国的一切尊荣。被剥夺了这一切恩惠，他们像裸体的逃亡者被驱逐出境。而在他们的位置上，卑微的人被举扬，即那些饥渴慕义的外邦民族。因为他们发现了自己的卑微和对天主知识的饥渴，恳求天主的话语，哪怕是碎屑，如同客纳罕妇人一般；因此他们充满了神圣奥秘的富饶。因为生于童贞女的基督，我们的天主，已将神圣恩惠的全部产业交给了外邦人。\\BibleQuote{祂扶助了祂的仆人以色列。} 不是一般的以色列，而是那真正保持以色列高贵身份的仆人。为此，天主之母也称为祂为仆人（子）和继承者。因为当祂发现那仆人在文字和法律中辛苦劳作时，祂以恩宠召叫了他。因此，这样的以色列，祂召叫并扶助了，以纪念祂的仁慈。\\BibleQuote{正如祂对我们的先祖所说的，对亚巴郎和他的子孙直到永远。} 这几句话包含了整个救恩工程的奥秘。因为为拯救人类并实践与先祖所立的盟约，基督曾\\BibleQuote{屈尊诸天并降下。} 于是，祂按我们所能接受的方式显示给我们，使我们能看见祂、触摸祂并听见祂说话。为此，天主圣言认为应藉一个女人——荣福童贞——取血肉和完全的人性，并出生为人，以偿付我们的债务，并在祂自己身上满全与亚巴郎所立盟约的规定，包括割损礼及与之相关的所有法律规条。说了这些话后，荣福童贞回到纳匝肋；此后，因凯撒的一道谕令，她再次来到白冷；因为她本人出自王室，便与她的未婚夫若瑟一同来到达味的王室。在那里发生了超乎一切奇事的奥秘——童贞女生子，并亲手抱着那以自己话语支撑整个受造界的那一位。\\BibleQuote{在客栈中为他们没有地方。} 那以自己话语建立全地的人，竟找不到地方。她用乳汁喂养那赐予一切有气息者生命和滋养的那一位。她用襁褓包裹那以自己话语束缚整个受造界的那一位。她将祂放在马槽中，那一位却坐在革鲁宾之上。天上的光环绕祂，那一位光照整个受造界。天军以赞颂侍奉那一位在万世之前已在诸天受光荣者。一颗星用它的火炬引导那些从远方来寻找那真东方的人。从东方来献礼者向那为我们而成为贫穷者献礼。荣福童贞保存了这些话，在自己心中默想，如同所有奥秘的容器。至圣童贞，你的赞美超越一切颂扬，因那由你取得血肉并出生为人的天主。天上、地上、地下的一切受造物都向你呈上当有的荣耀。因为你被显明为真正的革鲁宾之座。你在诸智王国的至高之处闪耀如光明的本身；那里，无始的父得光荣，祂的能力曾荫庇你；那里，圣子受朝拜，你按肉身生了祂；那里，圣神受赞颂，祂在你子宫中成就了大王的诞生。蒙恩的女子，藉着你，圣而共融的天主圣三在世上被认识。请使我们与你一起，堪当在耶稣基督我们的主内，分享你圆满的恩宠：愿光荣归于父、及子及圣神，起初如何，今日亦然，直到永远。亚孟。\par

\SourceRecord{Translated by S.D.F. Salmond.；Ante-Nicene Fathers；Vol. 6.；Edited by Alexander Roberts, James Donaldson, and A. Cleveland Coxe.；Buffalo, NY: Christian Literature Publishing Co.,；1886.；<http://www.newadvent.org/fathers/06092.htm>.}

% structure: p0001 source=h1 target=section reason=page-title

\section{第三篇讲道}

论向童贞玛利亚的领报\par

Again have we the glad tidings of joy, again the announcements of liberty, again the restoration, again the return, again the promise of gladness, again the release from slavery. An angel talks with the Virgin, in order that the serpent may no more have converse with the woman. In the sixth month, it is said, the angel \\Person{加俾额尔} was sent from God to a virgin espoused to a man. \\Person{加俾额尔} was sent to declare the world-wide salvation: \\Person{加俾额尔} was sent to bear to \\Person{亚当} the signature of his restoration; \\Person{加俾额尔} was sent to a virgin, in order to transform the dishonour of the female sex into honour; \\Person{加俾额尔} was sent to prepare the worthy chamber for the pure spouse; \\Person{加俾额尔} was sent to wed the creature with the Creator; \\Person{加俾额尔} was sent to the animate palace of the King of the angels; \\Person{加俾额尔} was sent to a virgin espoused to \\Person{若瑟}, but preserved for \\Person{耶稣} the Son of God. The incorporeal servant was sent to the virgin undefiled. One free from sin was sent to one that admitted no corruption. The light was sent that should announce the Sun of righteousness. The dawn was sent that should precede the light of the day. \\Person{加俾额尔} was sent to proclaim Him who is in the bosom of the Father, and who yet was to be in the arms of the mother. \\Person{加俾额尔} was sent to declare Him who is upon the throne, and yet also in the cavern. The subaltern was sent to utter aloud the mystery of the great King; the mystery, I mean, which is discerned by faith, and which cannot be searched out by officious curiosity; the mystery which is to be adored, not weighed; the mystery which is to be taken as a thing divine, and not measured. \\BibleQuote{在第六个月，\\Person{加俾额尔}被派遣到一位童贞女那里。} What is meant by this sixth month? What? It is the sixth month from the time when \\Person{依撒伯尔} received the glad tidings, from the time that she conceived \\Person{若翰}. And how is this made plain? The archangel himself gives us the interpretation, when he says to the virgin: \\BibleQuote{请看，你的亲戚\\Person{依撒伯尔}，她在老年也怀了一个男胎；她现在已是第六个月了，她原是被称为不生育的。} In the sixth month—that is evidently, therefore, the sixth month of the conception of John. For it was meet that the subaltern should go before; it was meet that the attendant should precede; it was meet that the herald of the Lord's coming should prepare the way for Him. In the sixth month the angel \\Person{加俾额尔} was sent to a virgin espoused to a man; espoused, not united; espoused, yet kept intact. And for what purpose was she espoused? In order that the spoiler might not learn the mystery prematurely. For that the King was to come by a virgin, was a fact known to the wicked one. For he too heard these words of \\BibleBook{}{依撒意亚}: \\BibleQuote{看，一位童贞女将怀孕，生子。} And on every occasion, consequently, he kept watch upon the virgin's words, in order that, whenever this mystery should be fulfilled, he might prepare her dishonour. Wherefore the Lord came by an espoused virgin, in order to elude the notice of the wicked one; for one who was espoused was pledged in fine to be her husband's. \\BibleQuote{在第六个月，天使\\Person{加俾额尔}被派遣到一位童贞女那里，她已许配给一个名叫\\Person{若瑟}的男人。} Hear what the prophet says about this man and the virgin: \\BibleQuote{这封上的书卷要交给一个识字的人。} What is meant by this sealed book, but just the virgin undefiled? From whom is this to be given? From the priests evidently. And to whom? To the artisan \\Person{若瑟}. As, then, the priests espoused \\Person{玛利亚} to \\Person{若瑟} as to a prudent husband, and committed her to his care in expectation of the time of marriage, and as it behooved him then on obtaining her to keep the virgin untouched, this was announced by the prophet long before, when he said: \\BibleQuote{这封上的书卷要交给一个识字的人。} And that man will say, I cannot read it. But why can you not read it, O \\Person{若瑟}? I cannot read it, he says, because the book is sealed. For whom, then, is it preserved? It is preserved as a place of sojourn for the Maker of the universe. But let us return to our immediate subject. In the sixth month \\Person{加俾额尔} was sent to a virgin—he who received, indeed, such injunctions as these: \\Quote{现在前来，总领天使，成为这隐藏的可怕奥迹的仆役，你要作这奇迹的代理人。我被我的怜悯所感动，要降临大地，以恢复失落的\\Person{亚当}。罪使我按自己肖像所造的人腐败，败坏了我手的工程，蒙蔽了我所塑造的美。狼吞吃我养育的，乐园的家园荒凉，生命树被火焰的剑把守，快乐的地方关闭了。我对这仇敌的对象生发怜悯，我要擒拿那敌人。但我愿将这奥迹——我仅向你透露——仍隐藏于所有天上的权力之外。因此，到童贞\\Person{玛利亚}那里去吧。你前往那有生命的城，关于它先知曾这样说：\\InnerQuote{光荣的事已论到你，天主的城。} 那么，前往我理性的乐园；前往东门；前往配得上我圣言的居所；前往那地上的第二个天；前往那轻云，宣布我来临的阵雨；前往为我预备的圣所；前往降生成人的厅堂；前往我按肉身出生的纯洁内室。在我理性方舟的耳中说话，为我预备听觉的通道。但不要扰动或烦扰童贞女的灵魂。以适合那圣所的方式显现，先以喜乐的声音问候她。并用这话向\\Person{玛利亚}致意：\\InnerQuote{万福！充满恩宠者！} 好使我怜悯\\Person{厄娃}的堕落。} The archangel heard these things, and considered them within himself, as was reasonable, and said: \\Quote{这事奇怪，所说的事超出理解。他是\\OriginalTerm{革鲁宾}{Cherubim}所畏惧的，\\OriginalTerm{色辣芬}{Seraphim}不能仰视的，一切天上的权力所不能理解的，他竟保证与一女子结合吗？他宣布自己亲自来临吗？不仅如此，他还提供听觉的通道吗？那定\\Person{厄娃}罪的，竟急于如此尊崇她的女儿吗？因为他曾说：\\InnerQuote{为我预备听觉的通道。} 但子宫怎能容纳那不能被空间所容的呢？这实在是可怕的奥迹。} While the angel is indulging such reflections, the Lord says to Him: \\Quote{你为什么烦恼困惑呢，\\Person{加俾额尔}？你不是已经被我派遣到司祭\\Person{匝加利亚}那里去了吗？你不是向他传达了\\Person{若翰}诞生的喜讯吗？你不是惩罚了那不信的司祭，使他哑口无言吗？你不是使他变哑了吗？你不是宣告了，而我也证实了吗？事实不也随着你的喜讯而发生了吗？那不生育的妇人不是怀孕了吗？子宫不是听从了话吗？不生育的病症不是消失了吗？自然惰性的倾向不是逃走了吗？她从前从未怀孕，如今不是显出多产了吗？对我，万物的创造者，有什么事是不可能的呢？那么，你为什么又摇摆不定呢？} What is the angel's answer to this? \\Quote{主啊，} he says, \\Quote{补救自然的缺陷，消除邪恶的冲击，召回死亡的肢体恢复生命的能力，命令自然有生育的效能，除去超过一般年限的肢体的不生育，将枯干的旧茎变成青翠活力的样子，使不结果的土地突然成为出产谷穗的地方——做这一切工作，正如通常的情形，需要你的大能。\\Person{撒辣}是见证，与她一起的还有\\Person{黎贝加}和\\Person{亚纳}，她们都曾被不生育的可怕疾病束缚，后来都蒙你恩赐脱离那病患。但一位童贞女不与人接触就生育，这超越了一切自然法则；而你竟向这少女宣布你的来临？天地界限不能容纳你，童贞女的子宫怎能容纳你呢？} And the Lord says: \\Quote{\\Person{亚巴郎}的帐幕如何容纳了我？} And the angel says: \\Quote{在那里因着深厚的款待，主啊，你在帐幕门口向\\Person{亚巴郎}显现，并迅速经过它，如充满万物的那一位。但\\Person{玛利亚}怎能承受神性的火？你的宝座闪耀着光辉的照明，童贞女能不焚毁而接你吗？} Then the Lord says: \\Quote{的确，如果旷野中的火伤害了荆棘，我的来临也会伤害\\Person{玛利亚}；但如果那作为神性之火从天上降临的预像的火，使荆棘茂盛而不烧毁它，那么，对于那不以火焰而是以雨的形式降临的真理，你将说什么呢？} Thereupon the angel set himself to carry out the commission given him, and repaired to the Virgin, and addressed her with a loud voice, saying: \\Quote{\\BibleQuote{万福！充满恩宠者！主与你同在。} 从此，魔鬼不再反对你；因为从前那仇敌在哪里施予创伤，现在医师便首先在那里敷上救赎的药膏。死亡从哪里出来，生命也在那里预备了入口。由女人来了我们祸患的洪流，也由女人来了我们福泽的泉源。万福！充满恩宠者！不要羞愧，好像你是我人定罪的原因；因为你已成了那同时是审判者和救赎者的母亲。万福！你这世界失偶新郎的无玷母亲！万福！你将（来自母亲\\Person{厄娃}的）死亡沉入你的子宫！万福！你是有生命的天主圣殿！万福！你是天地同等的家园！万福！你是无限本性的最宽广的容器！} But as these things are so, through her has come for the sick the Physician; for them that sit in darkness, the Sun of righteousness; for all that are tossed and tempest-beaten, the Anchor and the Port undisturbed by storm. For the servants in irreconcilable enmity has been born the Lord; and One has sojourned with us to be the bond of peace and the Redeemer of those led captive, and to be the peace for those involved in hostility. For \\BibleQuote{他是我们的和平}；愿我们众人均能获享这和平，赖我们的主耶稣基督的恩宠和慈爱；愿光荣、尊威和权能归于他，从今直到永远，世世无穷。阿们。\par

\SourceRecord{Translated by S.D.F. Salmond.；Ante-Nicene Fathers；Vol. 6.；Edited by Alexander Roberts, James Donaldson, and A. Cleveland Coxe.；Buffalo, NY: Christian Literature Publishing Co.,；1886.；<http://www.newadvent.org/fathers/06093.htm>.}

% structure: p0001 source=h1 target=section reason=page-title

\section{第四篇讲道词}

论神圣的显现，或论基督受洗\par

你们这些基督的朋友、客旅的朋友和弟兄的朋友，请和善地接纳我今天的讲话，敞开你们的耳朵如同听的门，让我的话语进入其中，并从我这里接受这关于基督圣洗的救恩宣告——它发生在约旦河中——为的是你们对主的爱慕之情可以因祂如此屈尊俯就为我们所做的一切而更加热切。因为即使救主主显节的庆典已经过去，祂的恩宠仍常与我们同在。因此，让我们以永不知足的心来享用它；因为对救恩之事而言，不知足的渴求是好的——是的，是好的。所以，我们众人都从加里肋亚到犹太去吧，让我们与基督一同前行；因为与这样的同伴在生命之路上旅行是有福的。来吧，让我们用思想的脚步前往约旦河，看看洗者若翰如何为那不需要圣洗却仍接受此礼的那一位施洗，为的是将圣洗的恩宠白白赐予我们。来吧，让我们在这水中所象征的图像中观看我们重生的景象。\BibleQuote{那时，耶稣从加里肋亚来到约旦河若翰那里，要受他的洗。}哦，主的谦卑何其浩大！哦，祂的屈尊俯就何其浩大！诸天的君王急忙来到若翰——祂自己的前驱——面前，没有调动祂的天使军团，没有预先派遣无形体的大能作为祂的前驱；而是以极致的简朴、以士兵的形态，来到祂自己的下属面前。祂走近他如同群众中的一员，在俘虏中自卑，却原是救赎者；与被审判者同列，却原是审判者；与亡羊为伍，却原是善牧——祂为迷途的羊从天上降下，却未曾离开祂的诸天；与莠子混杂，却原是那无需播种的天上谷粒。当洗者若翰看见祂时，认出祂——就是他从前在母腹中已经认出并朝拜的那一位——并且清楚意识到这就是那一位，为祂的缘故，他超越自然的时间，在母腹中跳跃，突破了本性的界限，便将右手缩进双层外套里，像充满爱意的仆人向主人低头，对祂说：\Quote{我需要受祢的洗，祢反而到我这里来？我的主，祢这是做什么？祢为什么要颠倒秩序？祢为什么与仆人一同，在祢的仆人手中寻求属于仆人的事物？祢为什么想要接受祢并不需要的东西？祢为什么要用这浩大的屈尊俯就来压迫我，祢的仆役？我需要受祢的洗，但祢不需要受我的洗。较小的蒙较大的祝福，较大的并不被较小的祝福和圣化。光明由太阳点燃，太阳并不被油灯照亮。泥土由陶匠塑造，陶匠并不被泥土塑造。受造物由造物主更新，造物主并不被受造物恢复。病弱者由医生医治，医生并不被病弱者治愈。穷人从富人那里得到资助，富人并不向穷人借贷。我需要受祢的洗，祢反而到我这里来？难道我不知道祢是谁，祢从哪里发出光明，祢从哪里来吗？或者，因为祢也如我一样出生，我就要否认祢天主性的伟大吗？或者，因为祢如此俯就我，以至于接近我的身体，将整个我承载在祢内，为成就全人的救恩，我就因那可见的身体而忽视那仅可领悟的天主性吗？或者，为了我的救恩，祢已将自己献上的初果，我就忽视祢\BibleQuote{以光作外衣}吗？或者，因为祢穿着与我相关的肉身，以人类所能看见的方式显现在他们面前，我就忘记祢荣耀天主性的光辉吗？或者，因为我看见祢身上我的形象，我就与祢那不可见、不可理解的天主性作对吗？主啊，我认识祢；我清楚地认识祢。我认识祢，因为祢教导了我；因为除非享受祢的光照，无人能认识祢。主啊，我清楚地认识祢；因为我尚未看见这光时就已以灵性看见祢。当祢完全在天上圣父的无形怀抱中时，祢也完全在祢婢女和母亲的子宫内；而我虽因本性被拘禁在依撒伯尔的子宫中如同监狱，被未出生的孩子不能解脱的捆绑所束缚，却以预先的欢喜跳跃庆祝祢的诞生。那么，我这位在祢出生前就已预告祢的尘世旅居的人，能否在祢出生后不理解祢的来临？我这位在母腹中就是祢来临的老师的人，现在在完美的知识面前难道要成为理解上的孩童吗？但我不能不朝拜祢——整个受造界都朝拜祢；我不能不宣告祢——诸天曾以星辰指明祢，大地曾以贤士热情接待祢，天使的众军也因祢对我们的屈尊俯就而欢喜赞美祢，夜间守更的牧羊人歌颂祢为理性羊群的最高牧者。在祢面前我不能保持沉默，因为我是声音；是的，正如经上所说，\BibleQuote{在旷野中呼喊者的声音：预备主的道路。}我需要受祢的洗，祢反而到我这里来？我出生，从而消除了生我母亲的荒胎；我还是婴儿时就医治了父亲的口哑，从小就从祢那里得了行奇事的恩赐。但祢，从童贞玛利亚出生，按祢的意愿，也只有祢知道，并没有消除她的童贞；祢反而保全了它，只赐予她\Emph{母亲}的名号。她的童贞既没有阻碍祢的出生，祢的出生也没有损害她的童贞。但这两种全然相反的事——生育与童贞——却因同一意向而和谐；因为这样的事，在祢——本性的创造者——那里是可能的。我是人，是神圣恩宠的分享者；但祢是神，也是人，同样：因为祢本性就是人的朋友。我需要受祢的洗，祢反而到我这里来？祢是太初就有的，与天主同在，就是天主；祢是圣父荣耀的光辉；祢是完美圣父的完美肖像；祢是真光，照亮每一个来到世界的人；祢在世界中，来到祢所在之处；祢成了血肉，却没有变为血肉；祢住在我们中间，以仆人的形象向祢的仆人显现；祢以祢的圣名将天地连接起来——祢到我这里来？如此伟大的一位到像我这样渺小的人这里来？君王到前驱这里来？主到仆人这里来？虽然祢不以卑微人性的尺度降生为耻，但我没有能力超越本性的尺度。我知道地与创造主之间的差距有多大。我知道泥土与陶匠之间的区别有多大。我知道祢——公义的太阳——比我——不过是祢恩宠的火把——具有多么大的超越性。即使祢被身体的纯净云彩所包围，我仍能认出祢的主权。我承认我的奴役，我宣告祢的荣耀伟大，我承认祢的完美主权，我承认我自己的完全微不足道；我不配给祢解鞋带，我怎敢触摸祢无罪的头顶？我怎能向祢伸出右手，祢铺张诸天如帐幕，将地立在众水之上？我怎能将这些仆人的手按在祢头上？我怎能洗祢——无玷无罪的祢？我怎能照亮光明？我该为祢献上怎样的祈祷？祢甚至接受那些不认识祢之人的祈祷？}\par

当我为别人授洗时，我因你的名授洗，为使他们信你——你将带着荣耀来临；但当我为你授洗时，我将提及谁？我将因谁的名为你授洗？是因父的名？但你完全在父内，父也完全在你内。或是因子的名？但除你之外，没有别的天主子是出于本性的。或是因圣神的名？但圣神永远与你同在，因为祂与你\OriginalTerm{同体}{consubstantial}、同旨、同判、同能、同荣；祂与你一同从万有接受同样的钦崇。因此，主啊，如果你愿意，请为我授洗；请为洗者授洗。请重生那由你所生的人。请伸出你可畏的右手，就是你为自己预备的那只手，以你的触摸加冕我的头，使我能在你的国度前奔跑，像一位先驱者戴冠，并殷勤地向罪人报告喜讯，以这热切的呼声向他们讲话：\Quote{请看，天主的羔羊，除免世罪者！}约旦河啊，请与我一同欢唱，与我一同跳跃，有节奏地搅动你的水流，如同舞蹈的律动；因为你的创造主正以肉身站在你旁边。昔日你曾看见以色列经过你，你分开你的洪流，等待百姓经过；但现在你更要决然地分开，更顺畅地流淌，拥抱祂那无瑕的肢体——祂曾在古时带领犹太人经过你。群山和丘陵，山谷和洪流，海洋和江河，请赞美上主，祂来到了约旦河上；因为通过这些水流，祂将圣化传递给一切水流。耶稣回答他说：\Quote{你暂且容许吧，因为我们应当这样，以完成全部义德。你暂且容许吧，洗者，请对\OriginalTerm{救恩经济}{economy}的时刻保持静默。学会愿意我所愿意的一切。学会在我所决心的事上服务我，不要好奇地探究我所要做的一切。你暂且容许吧：现在不要宣告我的神性；现在不要用你的嘴唇预报我的国度，免得暴君知道这事而放弃他对我的图谋。容许魔鬼来到我面前，与我交战，好像我只是一个普通人，并这样接受他的致命伤。容许我完成我来到地上的目的。今天在约旦河所进行的是一个\OriginalTerm{奥迹}{mystery}。我的奥迹是为我和我自己的。这里有一个奥迹，不是为满足我自己的需要，而是为给那些受伤者设计疗药。这里有一个奥迹，它在这水中给予人类重生之天上江河的表征。你暂且容许吧：当你看到我在我手所做的工中，以合乎神性的方式行我看为好的事时，那时你再以赞美配合已完成的行动。当你看到我洁净癞病人时，那时你再宣告我是自然的创造者。当你看到我使瘸子成为敏捷的奔跑者时，你也以更快的脚步准备你的舌头赞美我。当你看到我驱逐魔鬼时，那时你再以钦崇欢呼我的国度。当你看到我以我的话使死人从坟墓中复活时，那时你再与复活的人一同荣耀我为生命之主。当你看到我在父的右边时，那时你再承认我是神圣的，与父和圣神平等，同为宝座、永恒和尊荣。你暂且容许吧；因为我们应当这样，以完成全部义德。我是立法者，也是立法者之子；我应当首先经过所有已确立的，然后到处揭示我恩赐的预兆。我应当完成法律，然后赐予恩宠。我应当先出示影子，然后出示实体。我应当完成旧约，然后颁布新约，将它写在人的心版上，用我的血签署它，用我的圣神印证它。我应当上十字架，被钉子刺穿，按照那能受苦的本性受苦，并以我的受苦治愈痛苦，藉树木治愈由树木媒介所加于人的创伤。我应当下到坟墓的最深处，为那些被拘禁在那里的死者。我应当藉着我在肉体内三天的解体，摧毁古老仇敌死亡的力量。我应当为那些坐在黑暗和死影中的人点燃我身体的火炬。我应当以肉身升到我神性所在之处。我应当将那在我内统治的亚当引见给父。我应当完成这些事，因为为了这些事，我已站在我手所做的工之中。我应当暂时受这圣洗，然后赐予众人\OriginalTerm{同体}{consubstantial}圣三的圣洗。因此，洗者啊，请为这当前的救恩经济借给我你的右手，正如玛利亚为我的诞生借给我她的子宫。请将我浸入约旦河的水流中，正如生我的母亲用婴儿的襁褓包裹我。请赐予我你的圣洗，正如童贞女赐予我她的乳汁。请握住我的头，这头为色辣芬所敬畏。用你的右手握住这与你血缘相关的头。握住这头，自然使之可触摸。握住这头，正是为此目的由我和我父所形成。握住我的头，若有人虔诚地握住它，将使他免于遭遇船难。请为我授洗，我注定要用圣洗圣事、圣神和火为那些信我的人授洗：用圣洗，能洗去罪恶的污秽；用圣神，能使属地的成为属神的；用火，天生能烧尽过犯的荆棘。}洗者听了这些话后，便将心思转向救恩的对象，领悟了所领受的奥迹，执行了神圣的命令；因为他既虔诚又乐于服从。他缓慢地伸出右手，似乎既颤抖又欢喜，为主授了洗。那时，在场的犹太人，以及附近的、远处的，都彼此商议，这样对自己和彼此说：\Quote{难道我们毫无理由地认为若翰高于耶稣？难道我们毫无理由地认为前者大于后者？这圣洗本身不正证明洗者的优越吗？授洗者不是作为优越者，受洗者作为低劣者吗？}但他们因对救恩经济奥迹的无知而如此彼此胡言乱语时，唯独主、那\OriginalTerm{独生子}{Only-begotten}的本性之父，唯独那完美认识祂所无痛苦地生下的那一位，为了纠正犹太人错误的想象，开启了天闸，派遣圣神以鸽子的形式降下，落在耶稣的头上，以此指出新的诺厄，是的，诺厄的创造者，以及那遭遇船难的本性的好舵手。祂自己从天以清晰的声音呼叫说：\Quote{这是我的爱子，——就是那里耶稣，而不是若翰；受洗者，而不是授洗者；祂是在万世之前由我所生，而不是由匝加利亚所生；祂是玛利亚按肉身所生，而不是由依撒伯尔出乎意料所生；祂是童贞仍保持完整所结的果实，而不是从已消除的不育所发出的嫩枝；祂是曾与你们交往的那一位，而不是在旷野长大的那一位。这是我的爱子，我所喜悦的：我的儿子，与我自己同体，而非不同；按不可见的方面与我同体，按可见的方面与你们同体，却没有罪。这就是那与我一同造人的那位。这是我的爱子，我所喜悦的。我的这个儿子和玛利亚的这个儿子并非两个不同的人物；但这是我的爱子——这一个既是肉眼所见，也是理智所领悟的。这是我的爱子，我所喜悦的；你们要听祂。如果祂说：\InnerQuote{我与父原是一体}，你们要听祂。如果祂说：\InnerQuote{看见我的，就是看见父}，你们要听祂。如果祂说：\InnerQuote{派遣我的比我更大}，请使这声音适应救恩经济。如果祂说：\InnerQuote{人子是谁？}你们这样回答祂：}\par

\BibleQuote{你是基督，永生天主之子。}通过这些话，如同从圣父那里以雷霆形式从天上发出的，人类获得了光照：他们领悟了造物主与受造物之间、君王与士兵（臣仆）之间、工作者与工作之间的区别；并因信德得到坚固，他们借着若翰的洗礼，来到基督——我们的真天主——面前，祂\BibleQuote{以圣神及火施洗}。愿光荣归于祂、及圣父、及至圣而赋予生命的圣神，自今而始，及至永远，永世无尽。阿们。\par

\SourceRecord{Translated by S.D.F. Salmond.；Ante-Nicene Fathers；Vol. 6.；Edited by Alexander Roberts, James Donaldson, and A. Cleveland Coxe.；Buffalo, NY: Christian Literature Publishing Co.,；1886.；<http://www.newadvent.org/fathers/06094.htm>.}
